% 2026-11-23 W11 第 1 部。Michelman 399--403 與 Post 429--431 依 PDF 頁圖核過（PDF 頁減 396、減 427）。
% 399 頁註 1 之外，400 頁頁底另有一則未編號之註（講稿經過），正文稱「400 頁頁底未編號之註」。
% Post 註 6 作 "Id. at 430"，指 Michelman 一文；所引兩個定義實在 Michelman 400 頁（Post 註 5 亦引 400 頁），已於註腳說明。
% Bickel 原著未取得，只依 Michelman 399 頁對該語的轉述。Michelman 註 3 之判決歸類依案件題材，註 3 本身以 See, e.g. 概列。
\chapter{憲政主義與民主之衝突的問題設定}
\label{ch:10_problem}

\begin{qaquestion}
Michelman\ix{persons}{michelman@Michelman, Frank I.} 與 Post\ix{persons}{post@Post, Robert} 的文章同載於 86 Calif. L. Rev. 第 3 期，分別為 1996--97 年度 Brennan\ix{persons}{brennan@Brennan, William J., Jr.} Center 研討會的講座文章與評論。
\begin{enumerate}
\tightlist
\item 兩篇文章如何成稿，Post 如何看待憲法理論這項工作？
\item Michelman 如何定義憲政主義\ix{terms}{03@\zhlidxcat{憲政主義與憲法之效力基礎}!憲政主義}與民主\ix{terms}{01@\zhlidxcat{民主理論}!民主}，舉了哪三個問題，最高法院的終局解釋\ix{terms}{08@\zhlidxcat{憲法解釋與司法自制}!解釋}之主張與這組定義有何關係？
\item 這項衝突與 Bickel\ix{persons}{bickel@Bickel, Alexander M.} 的抗多數難題\ix{terms}{04@\zhlidxcat{違憲審查權限之歸屬}!對抗多數的難題}有何不同？
\item 問責民主觀\ix{terms}{01@\zhlidxcat{民主理論}!問責民主觀／福利民主觀}與自治\ix{terms}{01@\zhlidxcat{民主理論}!自治}觀、民粹論\ix{terms}{01@\zhlidxcat{民主理論}!民粹論}與個人論、制度難題\ix{terms}{01@\zhlidxcat{民主理論}!制度難題}各指什麼？
\end{enumerate}
\end{qaquestion}

\section{兩篇之成文經過}
\label{sec:10_problem-1}

Michelman 文 399 頁註 1 記載，Brennan Symposium Lectures 由 New York University School of Law 的 Brennan Center for Justice 主辦，經費來自 Berkeley 法學院 Thomas Jorde 教授的捐贈；致謝對象包括 Berkeley 與 N.Y.U. 兩所法學院的院長 Herma Hill Kay 與 John Sexton，四位評論人，以及就舊稿提供意見的 Larry Kramer\ix{persons}{kramer@Kramer, Larry D.} 與 Kenneth Winston。\footcite[399 n.1]{michelman1998}400 頁頁底另有一則未編號的註，交代講稿與刊本之間的差別：Michelman 於 1996 年 11 月 4 日在 Berkeley、1997 年 3 月 10 日在 N.Y.U. 各講一次，講稿事先送交評論人。Berkeley 場的評論人是 Post 與 Kathleen Sullivan\ix{persons}{sullivan@Sullivan, Kathleen M.}，N.Y.U. 場是 Ronald Dworkin\ix{persons}{dworkin@Dworkin, Ronald} 與 Don Herzog\ix{persons}{herzog@Herzog, Don}。因 Post 在 Berkeley、Dworkin 在 N.Y.U. 到場，Michelman 兩場的講稿略有不同，在 Berkeley 講「Post」的部分，在 N.Y.U. 講「Dworkin」的部分，刊行時合併改寫，並與四位評論人的評論一併發表。Michelman 原文第 III 部對 Post，第 IV 部對 Dworkin，兩者之間因此互有呼應；改寫時，凡評論所針對之處，Michelman 說「I have tried in the rewriting to stand by everything to which the trenchant comments were directed」〔改寫時力求保持評論所針對的一切〕。\footcite[400（頁底未編號之註）]{michelman1998}

Post 文 429 頁星號註記載，他任教於 Berkeley 法學院（Boalt Hall），並在 1978 年 10 月開庭期擔任 Brennan 大法官的 law clerk〔法官助理〕。\footcite[429 n.\textdagger]{post1998}正文開頭稱 Michelman 以民主與憲政主義之間的張力為題「appropriate」〔適切〕，並說這項張力自二十世紀初以來即為美國憲法理論的中心問題。\footcite[429]{post1998}

Post 對憲法理論的定位也在 429 頁。他寫道，憲法理論「primarily seeks to expose and clarify the principles immanent within the practice of constitutional adjudication」〔主要在揭示並釐清內在於憲法裁判實務的原則〕，所以不能把自己置於 Rawls\ix{persons}{rawls@Rawls, John} 的「veil of ignorance」〔無知之幕〕或 Habermas\ix{persons}{habermas@Habermas, Jürgen} 的「ideal speech situation」〔理想言說情境〕之後；憲法理論「is always, so to speak, within our tradition and our history; it is parasitic on the very practice it seeks to explain」〔總是處於我們的傳統與歷史之內，寄生於它所要說明的實務〕。\footcite[429 \& nn.1--2]{post1998}依此，Brennan 這類受尊崇的大法官之成就屬於「part of the data to be explained」〔有待說明的資料之一部分〕，Post 因而贊同 Michelman 把 Brennan 的工作置於探究中心。\footcite[429]{post1998}430 頁接著說，實務中的原則並非一致，成熟的實務充滿尚未解決的內在張力（註 4 引 MacIntyre\ix{persons}{macintyre@MacIntyre, Alasdair} 論活的傳統是「an historically extended, socially embodied argument」〔歷經歷史延展、體現於社會的論辯〕），憲法理論的任務因此是以理論上的清晰、廣度與抽象程度參與這些論辯，無意超越或結束它們。\footcite[429--430 \& n.4]{post1998}

\section{憲政主義與民主之定義}
\label{sec:10_problem-2}

Michelman 在 400 頁以「appears to mean something like this」〔看來意指大致如下〕引出兩個粗略的定義。憲政主義是「The containment of politics by a supervening law that stands beyond the reach of the politics it is meant to contain」〔以立於其所約束之政治所不及之處的上位法來圍堵政治〕，這種法律可稱為「law of lawmaking」〔造法之法\ix{terms}{03@\zhlidxcat{憲政主義與憲法之效力基礎}!造法之法}〕，「controls which further laws can be made and by what procedures」〔控制得制定何種進一步的法律、依何種程序〕。民主是「Popular political self-government」〔人民的政治自治〕，即「the people of a country deciding for themselves the contents of the laws that organize and regulate their political association」〔一國人民自行決定規整其政治結社之法律的內容〕。\footcite[400]{michelman1998}

兩個定義合在一起，推出的是衝突。依民主原則，人民應自行決定一切有良好道德與物質理由關切的「politically decidable matters」〔政治上可決定的事項〕；造法之法的內容正屬其中，而依憲政主義原則，造法之法須對民主權威設定界限。Michelman 稱兩項原則因此處於「relentless conflict」〔無休止的衝突〕。\footcite[400]{michelman1998}

Post 在 430 頁接受這兩個定義，稱 Michelman 對憲政主義的界定「accurately」〔準確〕，對民主的界定「convincingly」〔有說服力〕，並逐字引述兩段。\footnote{Post 430 頁註 5 標明憲政主義之定義引自 Michelman 400 頁；註 6 引民主之定義，作「Id. at 430」。「Popular political self-government」一句見 Michelman 一文 400 頁，430 頁是 Post 自己文章的頁碼，此處以 400 頁為準。}

\section{三個問題與終局解釋之主張}
\label{sec:10_problem-3}

Michelman 以三個「politically decidable questions」說明，每一個都問一國要不要有一部造法之法\ix{terms}{03@\zhlidxcat{憲政主義與憲法之效力基礎}!造法之法}：\footcite[400--401]{michelman1998}

\begin{longtable}{@{}L{0.30\linewidth}L{0.36\linewidth}L{0.28\linewidth}@{}}
\caption{Michelman 設問之三個造法之法及其註 3 所列判決（依案件題材歸類）。}\\
\toprule
議題 & 400--401 頁之問法 & 註 3 所列判決 \\
\midrule
\endfirsthead
\toprule
議題 & 400--401 頁之問法 & 註 3 所列判決 \\
\midrule
\endhead
種族 & 是否有一部造法之法，幾乎禁止政府在任何情況、基於任何理由，為任何考量種族的立法或行政 & Adarand Constructors, Inc. v. Pena\ix{cases}{02@美國!01@聯邦最高法院!1995Adarand Constructors, Inc. v. Pena（1995）@Adarand Constructors, Inc. v. Pena（1995）}, 515 U.S. 200 (1995)；Miller v. Johnson\ix{cases}{02@美國!01@聯邦最高法院!1995Miller v. Johnson（1995）@Miller v. Johnson（1995）}, 515 U.S. 900 (1995) \\
競選經費 & 是否有一部造法之法，嚴格限制政府規範政治競選中的金錢流動 & Buckley v. Valeo\ix{cases}{02@美國!01@聯邦最高法院!1976Buckley v. Valeo（1976）@Buckley v. Valeo（1976）}, 424 U.S. 1 (1976)；Colorado Republican Fed. Campaign Comm. v. FEC\ix{cases}{02@美國!01@聯邦最高法院!1996Colorado Republican Federal Campaign Committee v. FEC（1996）@Colorado Republican Federal Campaign Committee v. FEC（1996）}, 518 U.S. 604 (1996) \\
生育 & 是否有一部造法之法，嚴格限制政府規範人民關於性、受孕、繼續懷孕與為人父母之所為 & Roe v. Wade\ix{cases}{02@美國!01@聯邦最高法院!1973Roe v. Wade（1973）@Roe v. Wade（1973）}, 410 U.S. 113 (1973)；Griswold v. Connecticut\ix{cases}{02@美國!01@聯邦最高法院!1965Griswold v. Connecticut（1965）@Griswold v. Connecticut（1965）}, 381 U.S. 479 (1965)；Eisenstadt v. Baird\ix{cases}{02@美國!01@聯邦最高法院!1972Eisenstadt v. Baird（1972）@Eisenstadt v. Baird（1972）}, 405 U.S. 438 (1972)；Bowers v. Hardwick\ix{cases}{02@美國!01@聯邦最高法院!1986Bowers v. Hardwick（1986）@Bowers v. Hardwick（1986）}, 478 U.S. 186 (1986) \\
\bottomrule
\end{longtable}

註 3 以 See, e.g. 概列此八案，並未逐案分配給三個問題，表中的歸類依各案題材。\footcite[401 n.3]{michelman1998}

Michelman 的論證是：這些選擇對許多人在物質與道德上都極其重要，認真對待民主，它們似乎都在民主所及的範圍之內；它們同時是關於造法之法的選擇，而憲政主義的原則表明，至少有一部分造法之法的選擇須安穩地置於民主政治的決定權之外。事實上，這三組問題正是美國聯邦最高法院在解釋並適用美國既有的造法之法，亦即憲法時，現今為這個國家所作的決定。\footcite[401]{michelman1998}

Michelman 引兩個判決說明最高法院如何看待這些決定。Cooper v. Aaron\ix{cases}{02@美國!01@聯邦最高法院!1958Cooper v. Aaron（1958）@Cooper v. Aaron（1958）} 的九位大法官一致意見主要由 Brennan 執筆，宣告大法官「supreme in the exposition of the law of the Constitution」〔在闡釋憲法之法上至高〕，358 U.S. 1, 18 (1958)；Planned Parenthood of Southeastern Pa. v. Casey\ix{cases}{02@美國!01@聯邦最高法院!1992Planned Parenthood of Southeastern Pa. v. Casey（1992）@Planned Parenthood of Southeastern Pa. v. Casey（1992）} 的決定性複數意見語氣更強，自任「speak[ing] before all others for [the] constitutional ideals」〔在一切人之先為憲法理想\ix{terms}{03@\zhlidxcat{憲政主義與憲法之效力基礎}!憲法理想}發言〕，505 U.S. 833, 868 (1992)，所涉不止於法律，而及於憲法的理想。\footcite[401 nn.4--5]{michelman1998}註 5 並記 Dworkin 在 Life's Dominion 120、126 頁表示贊同。\footcite[401 n.5]{michelman1998}Michelman 說，這些宣示的作者「don't seem particularly worried about the people deciding for themselves the contents or even the spirit of the fundamental laws」〔似乎並不擔心由人民自己決定根本法\ix{terms}{03@\zhlidxcat{憲政主義與憲法之效力基礎}!根本法}的內容，乃至其精神〕，並問：因此須將他們視為民主的敵人嗎？他的回答是未必；依流傳的說法，視其詮釋賦予造法之法何種精神與內容，他們也可以是民主的朋友。\footcite[401]{michelman1998}這類說法有兩種變形，一種將民主理解為程序性的理想，一種理解為實質性的理想，在兩位評論人身上各有代表：Post 代表程序，Dworkin 代表實質。Michelman 的計畫是陪兩人各從其出發的一極走到兩者在悖論中會合之處（程序一極見第 \ref{ch:20_responsive} 部）。\footcite[401]{michelman1998}

\section{與抗多數難題之比較}
\label{sec:10_problem-4}

Michelman 在 399 頁以 Bickel\ix{persons}{bickel@Bickel, Alexander M.} 的名句「Counter-Majoritarian Difficulty」為起點，所描述的情形是「unelected judges deciding the legal validity of the enactments of popular assemblies and thereby effectively ruling the country」〔未經選舉的法官決定民意機構所制定法律的法律效力，因而實際統治國家〕。他自承是看得見這項難題、又認真對待民主、同時認為 Brennan 的職業生涯對國家貢獻極大的人，並說 Brennan 或許比任何其他美國法官都更足以代表這項難題。\footcite[399]{michelman1998}

400 頁說明本文的主要目的之一，是把憲政主義與民主之衝突的討論推到 Bickel 難題之外，因為這項衝突「has a depth and a poignancy that Bickel's delicately understated formulation fails to capture」〔具有 Bickel 精細而含蓄的表述所未能捕捉的深度與痛切〕。\footcite[400]{michelman1998}註 2 另列把這個現象稱為「government by judiciary」〔法官統治〕的著作，Berger 的 Government by Judiciary（1977）與 Boudin 的同名書（1932）。\footcite[400 n.2]{michelman1998}

399 頁對 Bickel 難題的描述以未經選舉的法官為主體，所問是法官裁判民意機構制定之法的效力。400 頁對衝突的陳述則從兩個定義出發：民主原則所涵蓋的決定事項包括造法之法\ix{terms}{03@\zhlidxcat{憲政主義與憲法之效力基礎}!造法之法}的內容，憲政主義原則要求造法之法立於政治所不及之處，兩項原則於是就這批事項互相排斥，衝突的成立不以任何特定機關的行為為前提。Michelman 在 403 頁說，制度難題\ix{terms}{01@\zhlidxcat{民主理論}!制度難題}並非 Brennan 這類法官的司法角色所造成，並反問這種角色是否反而是解答。\footcite[403]{michelman1998}依此，法官的角色在 Michelman 的設問裡是受檢驗的解答，難題另有來源，來源在下一節。

\section{民主之兩種理由與制度難題}
\label{sec:10_problem-5}

\subsection{問責民主觀與自治觀}

402 頁起，Michelman 問人們為何在乎政治安排是否「ostensibly democratic」〔表面上民主〕。一種理由是，人們希望不論由誰統治，統治都合乎被統治者的利益。大致的想法是：統治者之任職取決於民意多數的容忍，選舉與代表制度又能公平反映人口中利益的分布，政府便傾向於合宜地回應被統治者的利益。Michelman 稱之為「accountability」或「welfare」view〔問責民主觀\ix{terms}{01@\zhlidxcat{民主理論}!問責民主觀／福利民主觀}或福利民主觀〕。\footcite[402]{michelman1998}

Post、Dworkin 與 Michelman 自己則在乎民主的另一種理由：人民自己治理自己，「exercising their own charge over the politically decidable conditions of their lives」〔對其生活中政治上可決定的各種條件行使自己的主管〕，也就是人性尊嚴與自由中通稱「positive liberty」〔積極自由\ix{terms}{01@\zhlidxcat{民主理論}!積極自由}〕的那一面。Michelman 稱這種理解為「self-government」conception。依自治觀，政治安排若未能依某種民主理論所設想的方式服務於自治，即屬有瑕疵。註 6 與註 7 分別指向 Post 的 Constitutional Domains 273、278、283--284 頁與 Dworkin 的 Freedom's Law 21--22 頁；註 8 引同處為把積極自由當作自我決定而重視的根據。\footcite[402 nn.6--8]{michelman1998}Michelman 預告三人將在一種說明上會合：民主藉由使每個人有理由「identify his or her political agency with the lawmaking and other acts of political institutions, or to claim such acts as his or her own」〔將自己的政治能動性認同於政治機構的立法與其他行為，或主張這些行為為自己所有〕，來服務於自治。\footcite[402]{michelman1998}

\subsection{民粹論與個人論}

Michelman 依自治所指的行為主體，把自治觀再分為兩類。「populist」一類的主張是，終究在道德上要緊的是實現一國人民作為單一「self」的政治意志；Michelman 說，持此說者解釋多數決程序如何轉化為自治的政治集體時，只遇到「moderate difficulty」〔中等程度的困難〕。註 9 所舉為 Richard D. Parker\ix{persons}{parker@Parker, Richard D.} 的 Here the People Rule（1994）與 Akhil Reed Amar\ix{persons}{amar@Amar, Akhil Reed} 於 94 Colum. L. Rev. 457, 503 (1994) 的文章。\footcite[402--403 \& n.9]{michelman1998}

Post、Dworkin 與 Michelman 則承襲標準自由主義的觀點，認為終究在道德上要緊的是個人的 agency、自由、尊嚴與自治。三人都相信民族、人民與政治社群\ix{terms}{01@\zhlidxcat{民主理論}!社群}確實存在，且其歷史與命運值得關切；但 Michelman 不信，也不信任何評論人相信，一個國家、人民或政治社群有心智或五官，從而有「a capacity for consciously self-directive agency」〔有意識地自我指引的能動能力〕而使我們有終極的道德理由關切它。註 11 引 Post 在 Constitutional Domains 306 頁的話佐證：「Groups neither reason nor have an autonomous will; only persons do」〔群體不推理，也沒有自主意志，只有個人才有〕。\footcite[403 \& nn.10--11]{michelman1998}所以 Michelman 說一國的自治，所指是「everyone」，即逐一計算的每一個人之自治。\footcite[403]{michelman1998}

\subsection{制度難題}

在這個前提下，Michelman 提出他稱為制度難題\ix{terms}{01@\zhlidxcat{民主理論}!制度難題}（Institutional Difficulty）的問題。立法依其性質是制度性的亦即集體的事業；一個對某項制度事件持異議、且不可能憑自己的一票或其他政治行動左右結果的人，如何能視自己為透過該事件而自治？一個人如何能透過政治所制定、令他反感的法律而自治？\footcite[403]{michelman1998}Post 與 Dworkin 的回答是，民主對這個兼具個人與集體性質的自治之謎提供解答，即民主給予個人以理由上的依據（「a warrant in reason」），使其將自己的政治能動性認同於多數決制度的行為。\footcite[403]{michelman1998}

\section{Post 對問題設定之回應}
\label{sec:10_problem-6}

Post 接受上述定義後，在 430 頁指出 Michelman 的表述有一處「strange」〔奇怪〕。Michelman 提到某種「necessity」〔必然性〕，使人不可抑制地想把根本法的決定「hive off」〔劃出\ix{terms}{03@\zhlidxcat{憲政主義與憲法之效力基礎}!劃出}〕民主的程序性管轄之外，並說這種必然性使民主完全被擋在「the contents of a country's most basic laws, its laws of lawmaking」〔一國最基本的法律，即其造法之法\ix{terms}{03@\zhlidxcat{憲政主義與憲法之效力基礎}!造法之法}的內容〕這個它似乎迫切而正當地想進入的決定空間之外（Post 註 7 標示其出處為 Michelman 420 頁）。Post 的反駁是，憲法裁判實務中有明顯的典型事例使憲法意義服從人民意志\ix{terms}{01@\zhlidxcat{民主理論}!人民意志}的專斷指示，Michelman 的結論對它們未置一詞。\footcite[430 \& n.7]{post1998}

Post 設想美國憲法依 Article V 的程序，由國會兩院各三分之二提案、四分之三的州議會批准，增修條文廢止第二十六增修條文，明文規定未滿 21 歲者不得投聯邦公職選舉的票；一位 18 歲的公民依此向聯邦法院起訴，請求投票權。\footcite[430--431 \& n.8]{post1998}Post 稱此為「paradigmatically easy case」〔典型的易案\ix{terms}{08@\zhlidxcat{憲法解釋與司法自制}!易案}〕，引 Schauer\ix{persons}{schauer@Schauer, Frederick} 的 Easy Cases，說判原告勝訴的法官可說未能理解憲法裁判實務，不論她的個人信念為何；法官縱使認為新增修條文是錯誤而反民主的，仍會判原告敗訴，理由是增修條文依 Article V 的程序標準反映了人民意志。Post 的結論是：「popular sovereignty manifestly trumps constitutionalism」〔人民主權\ix{terms}{01@\zhlidxcat{民主理論}!人民主權}明顯勝過憲政主義〕。\footcite[431 \& n.9]{post1998}

Post 隨即承認這個設例「loaded the dice」〔在骰子上動手腳〕：增修條文是新近而明確的。Michelman 所關切的、也是多數憲法理論所關切的案件很不同，涉及「cruel and unusual punishments」「the equal protection of the laws」「due process of law」「the freedom of speech」這類年代久遠、語句概括而多義的條款。Post 同意 Michelman 的結論，認為美國憲法裁判的歷史使人難以把這類條款的解釋與適用視為承載人民意志的權威，即使這些條款當初是由人民制定；只要這些條款仍然提供理由，使普通立法所表現的民意歸於無效，Michelman 所指民主與憲政主義之間的張力就「rise insolubly to the fore」〔無解地浮現〕。\footcite[431 \& nn.10--11]{post1998}

圖~\ref{fig:W11-relations} 依本部各節，列 Michelman 與 Post 在問題設定上援引、區辨或反駁之理論及其刊本頁。

\begin{qafigure}[第 1 部所涉理論之關係]{Michelman 與 Post 之問題設定所援引、區辨或反駁之理論。頁碼為 Michelman、Post 二文刊本頁。\label{fig:W11-relations}}
% 第 1 部理論關係：Michelman 與 Post 之問題設定各自援引、區辨、反駁之理論。
% 內容依 Translation/Register/relations.tsv（W11-REL-001 至 011；012 MacIntyre 不入圖，節點限 7 個）。
% 色調：Michelman figA、Post figD，被引者 slate；卡片淡底加左側色條（skill diagram-design「動手前」第 2 條）。
\begin{tikzpicture}[qa,x=1mm,y=1mm,
    card/.style={fill=#1bg,text width=37mm,minimum height=12mm,inner xsep=3.5mm,inner ysep=2mm,anchor=north west,
      execute at begin node={\raggedright\setstretch{1.2}},
      path picture={\fill[#1] (path picture bounding box.north west) rectangle
        ([xshift=3pt]path picture bounding box.south west);}},
    rel/.style={figGray,line width=0.6pt,-{Stealth[length=4pt,width=3.2pt]},shorten >=1pt,shorten <=1pt},
    lab/.style={font=\scriptsize,text=dimink,inner sep=1pt,align=center,execute at begin node={\setstretch{1.15}}}]
  \node[card=figGray] (B) at (0,0) {Bickel\\{\scriptsize 對抗多數的難題}};
  \node[card=figGray] (RH) at (102,0) {Rawls、Habermas\\{\scriptsize 無知之幕、理想言說情境}};
  \node[card=figA] (M) at (0,-34) {Michelman\\{\scriptsize 憲政主義與民主之衝突}};
  \node[card=figD] (P) at (102,-34) {Post\\{\scriptsize 人民主權與憲法裁判實務}};
  \node[card=figGray] (PA) at (0,-70) {Parker、Amar\\{\scriptsize 民粹論}};
  \node[card=figGray] (D) at (51,-70) {Dworkin\\{\scriptsize 自治觀}};
  \node[card=figGray] (S) at (102,-70) {Schauer\\{\scriptsize 易案}};
  \draw[rel] (M.north) ++(-4,0) -- ([xshift=-4mm]B.south) node[lab,midway,right=1mm,align=left] {推到難題之外\\（399--400 頁）};
  \draw[rel] (M.south) ++(-4,0) -- ([xshift=-4mm]PA.north) node[lab,midway,right=1mm,align=left] {區辨，不採\\（402--403 頁註 9）};
  \draw[rel] ([xshift=-4mm]M.south east) -- (D.north) node[lab,pos=0.62,anchor=west,xshift=7mm,align=left] {同屬自治觀\\（402 頁註 7）};
  \draw[rel] ([yshift=3mm]M.east) -- ([yshift=3mm]P.west) node[lab,midway,above] {同屬自治觀（402 頁註 6）};
  \draw[rel] ([yshift=-3mm]P.west) -- ([yshift=-3mm]M.east) node[lab,midway,below=1mm] {接受兩個定義（430 頁）\\以修憲設例反駁其結論\\（430--431 頁）};
  \draw[rel] (P.north) ++(-4,0) -- ([xshift=-4mm]RH.south) node[lab,midway,right=1mm,align=left] {區辨：憲法理論\\不置於其後（429 頁）};
  \draw[rel] (P.south) ++(-4,0) -- ([xshift=-4mm]S.north) node[lab,midway,right=1mm,align=left] {援引易案\\（431 頁註 9）};
\end{tikzpicture}

\end{qafigure}
